\chapter*{Apresentação --- memorial}
\phantomsection\addcontentsline{toc}{chapter}{APRESENTAÇÃO --- MEMORIAL}
% \chapter* não atualiza a marca de cabeçalho (só a versão numerada o faz), e
% sem isto as páginas seguintes herdam o "SUMÁRIO" da seção anterior.
\markboth{Apresentação --- memorial}{Apresentação --- memorial}

\section*{Trajetória e formação}

Sou graduado em Ciências Econômicas pela Universidade Federal de Goiás
(UFG, 2024). Ainda durante a graduação, em 2023, ingressei no Laboratório de
Sensoriamento Remoto e Geoprocessamento (LAPIG/UFG) como bolsista, para atuar
no apoio à gestão de projetos. Aos poucos, além das atividades de gestão,
passei a integrar também atividades de pesquisa, fazendo parte da equipe do
projeto Pasto Legal (\textit{pasto.legal}) e realizando trabalhos de
desenvolvimento de automações, sites e ferramentas para o laboratório.

A convivência diária com um laboratório de referência em sensoriamento remoto
moldou minha formação complementar. Em 2025, realizei o Geocurso de
Classificação de Imagens com Google Earth Engine, ministrado no LAPIG pelo
pesquisador Bernard Oliveira. Em 2026, realizei o Geocurso Explorando
Embeddings Geoespaciais com Modelos Fundacionais (Clay e AlphaEarth),
ministrado pelo pesquisador Tiago Gonçalves Maia Geraldine. Ainda em 2026,
ministrei no LAPIG um curso de extensão sobre automação de fluxos de trabalho
com Google Apps Script, GitHub Actions e inteligência artificial, transferindo
para outros a prática que eu vinha construindo.

\section*{O percurso do mestrado}

Cheguei ao Programa de Pós-Graduação em Ciências Ambientais (PPGCIAMB) pela
relação estreita entre o LAPIG e o programa. O tema --- a dinâmica das
pastagens em Goiás e seus fatores socioeconômicos --- está no encontro
natural entre a minha formação de economista e a tradição do laboratório no
monitoramento do uso da terra, especialmente no mapeamento de pastagens.

A pesquisa começou com uma hipótese clássica: a pastagem como estágio
intermediário da ocupação (vegetação $\to$ pasto $\to$ lavoura). O que ela
se tornou foi decidido pelos dados. O percurso está documentado em 54
pipelines de análise, com índice lógico, narrativa de construção e um
registro de decisões metodológicas numeradas --- e foi esse percurso
exploratório, com suas confirmações e refutações, que fez emergir a tese da
``marcha ao norte'' apresentada neste texto. Esta dissertação nasce, assim,
da vontade de contar as histórias que a série riquíssima de quatro décadas do
MapBiomas guarda. Nesse caminho, os dados do Instituto Brasileiro de
Geografia e Estatística (IBGE) também foram fundamentais, tanto para
construir a narrativa central quanto para revelar faces que os mapas sozinhos
não mostram.

\section*{O método de trabalho e o uso de inteligência artificial}

Devo registrar com franqueza o meu ponto de partida: cheguei ao mestrado sem
experiência prática de programação --- apenas noções básicas de lógica, R e
Python --- e sem bagagem em geotecnologias. A ponte foi construída pela
inteligência artificial. A partir de 2026, com o apoio e o incentivo do
LAPIG, em especial do professor Laerte, comecei a integrar ferramentas de IA
aos fluxos da gestão financeira dos projetos do laboratório; à medida que
ganhei fluência, levei essa prática para o centro da pesquisa.

O fluxo de trabalho que desenvolvi pode ser descrito assim: a curiosidade
vem primeiro; dela nasce uma pergunta formulada com precisão; a IA me ajuda
a materializar os scripts necessários para buscar a resposta nos dados; e,
encontrado um achado, a IA me ajuda a materializar visualizações ---
inclusive interativas --- que o tornem legível e verificável. As perguntas,
as decisões metodológicas e a interpretação dos resultados permanecem
minhas; o que a IA me fornece é autonomia para programar no ritmo da
curiosidade.

Essa autonomia cobra uma contrapartida de rigor, e a pesquisa a paga de
forma auditável: um diário de 27 decisões metodológicas numeradas
(Apêndice~\ref{ap:decisoes}), auditorias sistemáticas que rastreiam cada
número exibido até o arquivo de dados que o gerou, e autocorreções
documentadas --- resultados que seriam publicados e foram superados por um
teste a mais. A transparência operacional desse arranjo está declarada na
metodologia (seção~\ref{sec:met-ia}); o que pertence a este memorial é o seu
sentido formativo: foi por esse caminho que a curiosidade de um economista
ganhou os instrumentos para interrogar quarenta anos de mapas.
